% Part: incompleteness
% Chapter: arithmetization-syntax
% Section: proofs-in-nd

\documentclass[../../../include/open-logic-section]{subfiles}

\begin{document}

\olfileid{inc}{art}{pnd}
\olsection{સ્વાભાવિક નિગમનમાં \usetoken{P}{derivation}}

\begin{explain}
!!{derivation}sનું અંકગણિતીકરણ કરવા
આપણે !!{derivation}sને સંખ્યાઓ વડે
નિરૂપવાં પડશે. !!{derivation}s એ
!!{formula}sનાં વૃક્ષો છે અને તેમાં
દરેક અનુમાનને એક કે બે ચિહ્નન હોય છે.
તેથી પુનરાવર્તી નિરૂપણ સ્વાભાવિક છે:
!!{derivation}ને એવી ટપલ તરીકે નિરૂપીએ
જેના ઘટકો છેલ્લા અનુમાનના આધારવાક્યો
તરફ જતા સીધા ઉપ-!!{derivation}sની સંખ્યા,
આ ઉપ-!!{derivation}sનાં નિરૂપણો,
અંતિમ !!{formula}, છેલ્લા અનુમાનનું
નિવૃત્તિ-ચિહ્નન અને છેલ્લા અનુમાનનો
પ્રકાર દર્શાવતો નંબર છે.
\end{explain}

\begin{defn}
  $\delta$ એ સ્વાભાવિક નિગમનમાં
  !!a{derivation} હોય, તો $\Gn{\delta}$ને
  નીચે પ્રમાણે આગમનાત્મક રીતે વ્યાખ્યાયિત કરીએ:
  \begin{enumerate}
  \item
    $\delta$ માત્ર ધારણા~$!A$નું બનેલું
    હોય, તો $\Gn{\delta}$ એ
    $\tuple{0, \Gn{!A}, n}$ છે.
    ધારણા !!{undischarged} હોય તો
    નંબર~$n$ એ~$0$ છે; અન્યથા તે
    ધારણાનું અંકીય ચિહ્નન છે.
  \item
    $\delta$નો અંત શૂન્ય, એક, બે કે ત્રણ
    આધારવાક્યોવાળા અનુમાનથી થતો હોય,
    તો અનુક્રમે $\Gn{\delta}$ આ છે:
    \begin{align*}
      & \tuple{0, \Gn{!A}, n, k}, \\
      & \tuple{1, \Gn{\delta_1}, \Gn{!A}, n, k}, \\
      & \tuple{2, \Gn{\delta_1}, \Gn{\delta_2}, \Gn{!A}, n, k}, \text{ અથવા}\\
      & \tuple{3, \Gn{\delta_1}, \Gn{\delta_2}, \Gn{\delta_3}, \Gn{!A}, n,
        k},
    \end{align*}
    અહીં $\delta_1$, $\delta_2$, $\delta_3$
    એ~$\delta$ના છેલ્લા અનુમાનનાં
    આધારવાક્યોમાં પૂરાં થતાં
    ઉપ-!!{derivation}s છે;
    $!A$ એ~$\delta$ના છેલ્લા અનુમાનનું
    ફલિતવાક્ય છે; $n$ એ છેલ્લા અનુમાનનું
    નિવૃત્તિ-ચિહ્નન છે (અનુમાન કોઈ ધારણા
    નિવૃત્ત ન કરતું હોય તો~$0$); અને
    છેલ્લા અનુમાનમાં વપરાયેલા નિયમ પ્રમાણે
    નીચેના કોષ્ટકમાંથી~$k$ મળે છે.
  
    \begin{tabular}{lccccccc}
      \text{નિયમ:} & \Intro{\land} & \Elim{\land} & \Intro{\lor} & \Elim{\lor} \\
      $k$: & 1 & 2 & 3 & 4  \\[1.5ex]
      \text{નિયમ:} &  \Intro{\lif} & \Elim{\lif} & \Intro{\lnot} & \Elim{\lnot} \\
      $k$: & 5 & 6 & 7 & 8 \\[1.5ex]
      \text{નિયમ:}  &  \FalseInt & \FalseCl & \Intro{\lforall} & \Elim{\lforall} \\
      $k$: & 9 & 10 & 11 & 12 \\[1.5ex]
      \text{નિયમ:} & \Intro{\lexists} & \Elim{\lexists} & \Intro{\eq} & \Elim{\eq} \\
      $k$: & 13 & 14 & 15 & 16 
    \end{tabular}
  \end{enumerate}
\end{defn}

\begin{ex}
  આ અત્યંત સરળ !!{derivation} જુઓ:
  \begin{prooftree}
    \AxiomC{$\Discharge{!A \land !B}{1}$}
    \RightLabel{\Elim{\land}}
    \UnaryInfC{$!A$}
    \DischargeRule{\Intro{\lif}}{1}
    \UnaryInfC{$(!A \land !B) \lif !A$}
  \end{prooftree}
  ધારણાની ગડેલ-સંખ્યા
  $d_0 = \tuple{0,
    \Gn{(!A \land !B)}, 1}$ થાય.
  $\Elim{\land}$ના ફલિતવાક્યમાં
  પૂરી થતી !!{derivation}ની ગડેલ-સંખ્યા
  $d_1 = \tuple{1,
     d_0, \Gn{!A}, 0, 2}$ થાય
  ($1$ એટલા માટે કે $\Elim{\land}$ને
  એક આધારવાક્ય છે, ફલિતવાક્ય~$!A$ની
  ગડેલ-સંખ્યા લેવાઈ છે, $0$ એટલા માટે
  કે કોઈ ધારણા નિવૃત્ત થતી નથી અને
  $2$ એ $\Elim{\land}$નો સંકેત નંબર
  છે). આખી !!{derivation}ની ગડેલ-સંખ્યા
  તેથી
  $\tuple{1, d_1, \Gn{((!A \land !B) \lif
      !A)}, 1, 5}$, એટલે કે
  \[
  \tuple{1, \tuple{1, \tuple{0, \Gn{(!A \land !B)}, 1}, \Gn{!A}, 0, 2},
    \Gn{((!A \land !B) \lif !A)}, 1, 5}.
  \]
  \sourcecorrection{OLINC-011}{સ્થિર અંગ્રેજી
  મૂળમાં ધારણાના સંકેત~$d_0$માં
  $\Gn{!A \land !B}$ છે, જ્યારે
  એ જ ધારણાની વિસ્તૃત ટપલમાં
  $\Gn{(!A \land !B)}$ છે. વાક્યરચનાની
  એકરૂપ સંકેતશ્રેણી માટે અહીં બંને
  જગ્યાએ બીજા સ્વરૂપને રાખ્યું છે.}
\end{ex}

\begin{explain}
!!{derivation}sનું નિરૂપણ નક્કી કર્યા પછી
આવા !!{derivation}sની ગડેલ-સંખ્યાઓ પર
આદિમ પુનરાવર્તી રીતે ક્રિયા કરી શકાય
અને તેમના આવશ્યક ગુણધર્મો તથા સંબંધો
વ્યક્ત કરી શકાય તે પણ બતાવવું પડે.
કેટલીક ક્રિયાઓ સરળ છે: દાખલા તરીકે,
!!a{derivation}ની ગડેલ-સંખ્યા~$d$ હોય,
તો $\fn{EndFmla}(d) = (d)_{(d)_0+1}$
તેના અંતિમ !!{formula}ની ગડેલ-સંખ્યા
આપે છે; $\fn{DischargeLabel}(d) = (d)_{(d)_0 + 2}$
નિવૃત્તિ-ચિહ્નન આપે છે અને બિનધારણા
કિસ્સામાં
$\fn{LastRule}(d) = (d)_{(d)_0 + 3}$
છેલ્લા અનુમાનનો પ્રકાર દર્શાવતો નંબર
આપે છે. માત્ર ધારણા હોય ત્યારે આ
વિધેયને શૂન્ય મૂલ્ય આપીને સર્વત્ર
વ્યાખ્યાયિત કરી શકીએ.
\sourcecorrection{OLINC-012}{સ્થિર અંગ્રેજી
મૂળમાં ત્રણ-ઘટક ધારણા-ટપલ માટે
$\fn{LastRule}(d)$નું દર્શાવેલું
ઘટક અસ્તિત્વમાં નથી. આ સમીકરણ
બિનધારણા કિસ્સે લાગુ પડે છે;
ધારણાના કિસ્સામાં શૂન્ય મૂલ્ય
મૂકવાથી આગળનો ગુણધર્મ સર્વત્ર
અર્થપૂર્ણ થાય છે.}
બીજી કેટલીક ક્રિયાઓ ઘણી કઠિન છે.
તેમને કેવી રીતે સંભાળવી તેનો ઓછામાં
ઓછો રૂપરેખાત્મક ખ્યાલ આપીશું. હેતુ
એ બતાવવાનો છે કે ``$\delta$ એ
$\Gamma$માંથી~$!A$ની
!!a{derivation} છે'' એ સંબંધ
$\delta$ અને~$!A$ની ગડેલ-સંખ્યાઓનો
આદિમ પુનરાવર્તી સંબંધ છે.
\end{explain}

\begin{prop}
  નીચેના સંબંધો આદિમ પુનરાવર્તી છે:
  \begin{enumerate}
  \item $!A$ એ~$\delta$માં
    ચિહ્નન~$n$વાળી ધારણા તરીકે આવે છે.
  \item $\delta$માં ચિહ્નન~$n$વાળી
    બધી ધારણાઓ $!A$ સ્વરૂપની છે
    (એટલે $\delta$માં ચિહ્નન~$n$
    વાપરીને ધારણા~$!A$ને
    !!{discharge} કરી શકીએ).
  \end{enumerate}
\end{prop}

\begin{proof}
  !!{formula}sની ગડેલ-સંખ્યાઓ અને
  !!{derivation}sની ગડેલ-સંખ્યાઓ વચ્ચેના
  અનુરૂપ સંબંધો આદિમ પુનરાવર્તી છે
  તે બતાવવાનું છે.
  \begin{enumerate}
  \item $\fn{Assum}(x, d, n)$ આદિમ
    પુનરાવર્તી છે તે બતાવવું છે. તે
    ત્યારે સત્ય છે જ્યારે $x$ એ
    ગડેલ-સંખ્યા~$d$ ધરાવતી
    !!{derivation}માં ચિહ્નન~$n$વાળી
    ધારણાની ગડેલ-સંખ્યા હોય.
    ગડેલ-સંખ્યા~$\tuple{0, x,
      n}$ ધરાવતી !!{derivation} એ~$d$ની
    ઉપ-!!{derivation} હોય ત્યારે આવું બને.
    !!{derivation}sનું આપેલું
    સંકેતબદ્ધીકરણ
    \olref[cmp][rec][tre]{sec}માં રજૂ થયેલા
    વૃક્ષોના સંકેતબદ્ધીકરણનો વિશિષ્ટ
    કિસ્સો છે. તેથી આદિમ પુનરાવર્તી
    વિધેય~$\fn{SubtreeSeq}(d)$ આ~$d$ની
    બધી ઉપ-!!{derivation}sની ગડેલ-સંખ્યાઓની
    શ્રેણી આપે છે (તેની લંબાઈ વધુમાં
    વધુ~$d$ છે). આમ વ્યાખ્યાયિત કરીએ:
    \[
    \fn{Assum}(x, d, n) \defiff \bexists{i<d}{(\fn{SubtreeSeq}(d))_i =
      \tuple{0, x, n}}.
    \]
  \item $\fn{Discharge}(x, d, n)$
    સંબંધ આદિમ પુનરાવર્તી છે તે બતાવવું છે.
    તે ત્યારે સત્ય છે જ્યારે ગડેલ-સંખ્યા~$d$
    ધરાવતી !!{derivation}માં ચિહ્નન~$n$
    ધરાવતી બધી ધારણાઓ ગડેલ-સંખ્યા~$x$
    ધરાવતા !!{formula} હોય. આ સંબંધ
    ત્યારે અને માત્ર ત્યારે સત્ય છે જ્યારે
    $\bforall{y<d}{(\fn{Assum}(y, d, n) \lif y
      = x)}$.
  \end{enumerate}
\end{proof}

\begin{prop}
  \ollabel{prop:followsby}
  ગડેલ-સંખ્યા~$d$ ધરાવતી
  !!{derivation}~$\delta$નું છેલ્લું
  અનુમાન યોગ્ય હોય ત્યારે અને માત્ર
  ત્યારે સત્ય થતો ગુણધર્મ
  $\fn{Correct}(d)$ આદિમ પુનરાવર્તી છે.
\end{prop}

\begin{proof}
  દરેક અનુમાન-નિયમ~$R$ માટે
  $\fn{FollowsBy}_R(d)$ સંબંધ આદિમ
  પુનરાવર્તી છે તે બતાવવાનું છે.
  $\fn{FollowsBy}_R(d)$ ત્યારે અને માત્ર
  ત્યારે સત્ય છે જ્યારે $d$ એ
  !!{derivation}~$\delta$ની ગડેલ-સંખ્યા
  હોય અને~$\delta$નું અંતિમ !!{formula}
  $\delta$ના સીધા ઉપ-!!{derivation}sમાંથી
  $R$ના યોગ્ય ઉપયોગ વડે ફલિત થતું હોય.

  એક સરળ કિસ્સો \Intro{\land} નિયમનો
  છે. $\delta$નો અંત યોગ્ય
  $\Intro{\land}$ અનુમાનથી થતો હોય,
  તો તે આ પ્રમાણે દેખાય:
  \begin{prooftree}
    \AxiomC{}
    \RightLabel{$\delta_1$}
    \DeduceC{$!A$}
    
    \AxiomC{}
    \RightLabel{$\delta_2$}
    \DeduceC{$!B$}

    \RightLabel{\Intro\land}
    \BinaryInfC{$!A \land !B$}
  \end{prooftree}
  ત્યારે $\delta$ની ગડેલ-સંખ્યા~$d$
  એ $\tuple{2, d_1, d_2,
    \Gn{(!A \land !B)}, 0, k}$ છે, જ્યાં
  $\fn{EndFmla}(d_1) = \Gn{!A}$,
  $\fn{EndFmla}(d_2) = \Gn{!B}$,
  $n=0$ અને $k=1$. તેથી
  $\fn{FollowsBy}_{\Intro\land}(d)$ને
  આ પ્રમાણે વ્યાખ્યાયિત કરીએ:
  \begin{multline*}
    (d)_0 = 2 \land \fn{DischargeLabel}(d) = 0 \land \fn{LastRule}(d) = 1 \land {}\\
    \fn{EndFmla}(d) = {}\\
    \Gn{(} \concat \fn{EndFmla}((d)_1) \concat \Gn{\land}
    \concat \fn{EndFmla}((d)_2) \concat \Gn{)}.
  \end{multline*}

  બીજું સરળ ઉદાહરણ $\Intro\eq$ નિયમ છે.
  તેને આધારવાક્યો નથી, એટલે ધારણાની
  જેમ $(d)_0 = 0$. તેને
  નિવૃત્તિ-ચિહ્નન પણ નથી, એટલે
  $n=0$. જોકે $!A$ કોઈ બંધ
  પદ~$t$ માટે $\eq[t][t]$ સ્વરૂપનું
  હોવું જોઈએ. અહીં આદિમ પુનરાવર્તી
  વ્યાખ્યા આ છે:
  \begin{multline*}
    (d)_0 = 0 \land \fn{DischargeLabel}(d) = 0 \land
    \fn{LastRule}(d) = 15 \land {}\\
    \bexists{t<d}{(\fn{ClTerm}(t) \land 
      \fn{EndFmla}(d) = {}\\
      \Gn{{\eq}(} \concat t \concat \Gn{,} \concat t \concat \Gn{)})}.
  \end{multline*}
  \sourcecorrection{OLINC-013}{સ્થિર અંગ્રેજી
  મૂળમાં આ નિયમના સૂત્રમાં
  છેલ્લો નિયમ~$15$ છે તે શરત
  છૂટી છે. વ્યાખ્યા જણાવેલા
  $\Intro\eq$ નિયમને જ ઓળખે
  તે માટે અહીં આ શરત ઉમેરી છે.}

  વધુ જટિલ ઉદાહરણમાં
  $\fn{FollowsBy}_{\Intro{\lif}}(d)$
  ત્યારે અને માત્ર ત્યારે સત્ય છે જ્યારે
  $\delta$નું અંતિમ !!{formula}
  $(!A \lif !B)$ સ્વરૂપનું હોય,
  $\delta_1$નું અંતિમ !!{formula}~$!B$
  હોય અને~$\delta$માં ચિહ્નન~$n$
  ધરાવતી દરેક ધારણા $!A$ સ્વરૂપની હોય.
  આને આદિમ પુનરાવર્તી રીતે વ્યક્ત કરીએ:
  \begin{multline*}
    (d)_0 = 1 \land \fn{LastRule}(d) = 5 \land {}\\
    \bexists{a<d}{(\fn{Discharge}(a, (d)_1, \fn{DischargeLabel}(d)) \land {}}\\
      \fn{EndFmla}(d) = (\Gn{(} \concat a \concat \Gn{\lif}
      \concat \fn{EndFmla}((d)_1) \concat \Gn{)}))
  \end{multline*}
  (અહીં $a$ને~$!A$ની ગડેલ-સંખ્યા
  તરીકે વિચારો).
  \sourcecorrection{OLINC-014}{સ્થિર અંગ્રેજી
  મૂળના આ સૂત્રમાં છેલ્લો
  નિયમ~$5$ છે તે શરત છૂટી છે;
  નિયમ ઓળખવા માટે અહીં ઉમેરી છે.}

  હવે \Intro{\lexists}નું ઉદાહરણ લો.
  અહીં~$\delta$નું છેલ્લું અનુમાન
  ત્યારે અને માત્ર ત્યારે યોગ્ય છે
  જ્યારે એવું !!a{formula}~$!A$,
  બંધ પદ~$t$ અને !!a{variable}~$x$
  હોય કે $\Subst{!A}{t}{x}$ એ
  !!{derivation}~$\delta_1$નું અંતિમ
  !!{formula} હોય અને
  $\lexists[x][!A]$ છેલ્લા અનુમાનનું
  ફલિતવાક્ય હોય. તેથી
  $\fn{FollowsBy}_{\Intro{\lexists}}(d)$
  ત્યારે અને માત્ર ત્યારે સત્ય છે જ્યારે
  \begin{multline*}
    (d)_0 = 1 \land \fn{DischargeLabel}(d) = 0 \land
    \fn{LastRule}(d) = 13 \land {} \\
    \bexists{a < d}{\bexists{x<d}{\bexists{t<d}{
          (\fn{ClTerm}(t) \land \fn{Var}(x)  \land {}}}}\\
    \fn{Subst}(a,t,x) = \fn{EndFmla}((d)_1) \land
    \fn{EndFmla}(d) = (\Gn{\lexists} \concat x \concat a)).
  \end{multline*}
  \sourcecorrection{OLINC-015}{સ્થિર અંગ્રેજી
  મૂળના આ સૂત્રમાં છેલ્લો
  નિયમ~$13$ છે તે શરત છૂટી છે;
  અહીં એ સ્પષ્ટ કરી છે.}

  હવે $\fn{Correct}(d)$ને આ પ્રમાણે
  વ્યાખ્યાયિત કરીએ:
  \begin{multline*}
    \fn{Sent}(\fn{EndFmla}(d)) \land \bigl[{}\\
    (\fn{LastRule}(d) = 1 \land
    \fn{FollowsBy}_{\Intro\land}(d)) \lor \dots \lor {}\\
    (\fn{LastRule}(d) = 16 \land \fn{FollowsBy}_{\Elim\eq}(d)) \lor {}\\
    \bexists{n<d}{\bexists{x<d}{(d = \tuple{0, x, n})}}\bigr].
  \end{multline*}
  \sourcecorrection{OLINC-016}{સ્થિર અંગ્રેજી
  મૂળના સૂત્રમાં છેલ્લાં વિકલ્પોની
  આસપાસ કૌંસ નથી. સામાન્ય
  સંયોજક-પ્રાથમિકતા પ્રમાણે
  $\fn{Sent}$ની શરત ફક્ત પહેલા
  નિયમને લાગુ પડે તેમ બને છે.
  અહીં આખા વિકલ્પ-સમૂહને કૌંસમાં
  રાખ્યો છે જેથી દરેક કિસ્સામાં
  અંતિમ સૂત્ર વાક્ય હોવું જરૂરી બને.}
  પહેલી પંક્તિથી~$d$નું અંતિમ
  !!{formula} વાક્ય છે તે નિશ્ચિત
  થાય છે. છેલ્લી પંક્તિ~$d$
  માત્ર ધારણા હોય તે કિસ્સો આવરી લે છે.
\end{proof}

\begin{prob}
  \olref[inc][art][pnd]{prop:followsby}ની
  જેમ નીચેના ગુણધર્મો વ્યાખ્યાયિત કરો:
  \begin{enumerate}
  \item $\fn{FollowsBy}_{\Elim{\lif}}(d)$,
  \item $\fn{FollowsBy}_{\Elim{\eq}}(d)$,
  \item $\fn{FollowsBy}_{\Elim{\lor}}(d)$,
  \item $\fn{FollowsBy}_{\Intro{\lforall}}(d)$.
  \end{enumerate}
  છેલ્લા માટે એ પણ બતાવવું પડશે કે
  ગડેલ-સંખ્યા~$d$ ધરાવતી
  !!{derivation}નું છેલ્લું અનુમાન
  વિશિષ્ટ ચલની શરત સંતોષે છે કે નહીં
  તે આદિમ પુનરાવર્તી રીતે ચકાસી શકાય.
  એટલે $\Intro{\lforall}$ અનુમાનનો
  વિશિષ્ટ ચલ~$a$ ન તો~$d$ના
  અંતિમ !!{formula}માં આવે, ન તો~$d$ની
  ખુલ્લી ધારણામાં આવે. આ માટે
  \olref[inc][art][pnd]{prop:openassum}નો
  આદિમ પુનરાવર્તી વિધાન
  $\fn{OpenAssum}$ વાપરી શકો.
\end{prob}

\begin{prop}
  \ollabel{prop:deriv}
  યોગ્ય !!{derivation}~$\delta$ની
  ગડેલ-સંખ્યા~$d$ હોય ત્યારે સત્ય
  થતો સંબંધ~$\fn{Deriv}(d)$ આદિમ
  પુનરાવર્તી છે.
\end{prop}

\begin{proof}
  !!^a{derivation}~$\delta$ ત્યારે યોગ્ય
  છે જ્યારે તેનું દરેક અનુમાન કોઈ
  નિયમનો યોગ્ય ઉપયોગ હોય, એટલે કે
  તેની દરેક ઉપ-!!{derivation}નો અંત
  યોગ્ય અનુમાનથી થતો હોય. તેથી
  $\fn{Deriv}(d)$ ત્યારે અને માત્ર
  ત્યારે સત્ય છે જ્યારે
  \[
  \bforall{i<\len{\fn{SubtreeSeq}(d)}}{\fn{Correct}((\fn{SubtreeSeq}(d))_i)}
  \]
\end{proof}

\begin{prop}
  \ollabel{prop:openassum}
  ગડેલ-સંખ્યા~$d$ ધરાવતી
  !!{derivation}~$\delta$ની
  !!a{undischarged} ધારણા~$!A$ની
  ગડેલ-સંખ્યા~$z$ હોય ત્યારે સત્ય
  થતો સંબંધ~$\fn{OpenAssum}(z, d)$
  આદિમ પુનરાવર્તી છે.
\end{prop}

\begin{proof}
  ધારણાની કોઈ ઘટના~$\delta$ની એવી
  ઉપ-!!{derivation}માં ચિહ્નન~$n$
  સાથે આવે જેનો અંત નિવૃત્તિ-ચિહ્નન~$n$
  ધરાવતા નિયમથી થાય, તો તે ઘટના
  !!{discharged} કહેવાય. તેથી~$!A$ની
  ઓછામાં ઓછી એક ઘટના~$\delta$માં
  !!{discharged} ન હોય તો તે
  $\delta$ની !!a{undischarged} ધારણા છે.
  સાવચેતી જરૂરી છે: $\delta$માં~$!A$ની
  !!{discharged} અને
  !!{undischarged} બંને પ્રકારની
  ઘટનાઓ હોઈ શકે.

  $\delta_0$, \dots, $\delta_k$ એવી
  શ્રેણી લો જેમાં $\delta_0 =
  \delta$, $\delta_k$ એ ધારણા
  $\Discharge{!A}{n}$ છે (કોઈ~$n$
  માટે) અને $\delta_{i+1}$ એ
  $\delta_i$ની સીધી ઉપ-!!{derivation}
  છે. ધારણાનું ચિહ્નન શૂન્ય હોય અથવા એવી શ્રેણી
  હોય કે તેના કોઈ $\delta_i$નો અંત
  નિવૃત્તિ-ચિહ્નન~$n$ ધરાવતા
  અનુમાનથી ન થતો હોય, તો~$!A$ એ
  $\delta$ની !!a{undischarged}
  ધારણા છે.

  આદિમ પુનરાવર્તી
  વિધેય~$\fn{SubtreeSeq}(d)$ આપણને
  $\delta$ની બધી ઉપ-!!{derivation}sની
  ગડેલ-સંખ્યાઓની શ્રેણી આપે છે.
  $\delta$ની ઉપ-!!{derivation}sના
  સંકેતોની કોઈ માર્ગ-શ્રેણીના બધા
  ઘટકો તેમાં આવે છે. ઘટક-સમાવેશ
  આદિમ પુનરાવર્તી સંબંધ છે:
  $\fn{Subseq}(s, s')$ ત્યારે અને
  માત્ર ત્યારે સત્ય છે જ્યારે
  $\bforall{i<\len{s}}{\lexists[j<\len{s'}][(s)_i = (s')_j]}$.
  સીધી ઉપ-!!{derivation} હોવું પણ
  આદિમ પુનરાવર્તી છે:
  $\fn{Subderiv}(d, d')$ ત્યારે અને
  માત્ર ત્યારે, જ્યારે
  $\bexists{j<(d')_0}{d = (d')_{j+1}}$.
  અંતિમ વૃક્ષમાં માર્ગની લંબાઈ અને
  તેના સભ્યોના સંકેતો મૂળ કોડથી
  મર્યાદિત છે; તેથી અગાઉનું
  $\fn{sequenceBound}$ વિધેય
  વાપરીને $B(d)=\fn{sequenceBound}(d,d+1)+1$
  બધા આવા માર્ગોના સંકેતોની
  આદિમ પુનરાવર્તી મર્યાદા આપે છે.
  આમ $\fn{OpenAssum}(z, d)$ને
  આ પ્રમાણે વ્યાખ્યાયિત કરીએ:
  \begin{multline*}
    \bexists{s<B(d)}{(\fn{Subseq}(s, \fn{SubtreeSeq}(d))
      \land (s)_0 = d \land {}} \\
    \bexists{n<d}{((s)_{\len{s} \tsub 1} = \tuple{0, z, n} \land {}}\\
      \bforall{i<(\len{s} \tsub 1)}{(\fn{Subderiv}((s)_{i+1}, (s)_i) \land {}}\\
      (n=0 \lor \fn{DischargeLabel}((s)_i) \neq n)))).
  \end{multline*}
  \sourcecorrection{OLINC-017}{સ્થિર અંગ્રેજી
  મૂળના માર્ગ-પરીક્ષણમાં ચાર સમસ્યાઓ
  છે. $\fn{Subseq}$નું આપેલું સૂત્ર
  ક્રમયુક્ત ઉપશ્રેણી નહીં, માત્ર
  ઘટક-સમાવેશ તપાસે છે; માર્ગનો ક્રમ
  આગળની સીધી-ઉપપ્રમાણ શરત આપે છે.
  સીધી ઉપપ્રમાણની ટપલનું પ્રથમ સ્થાન
  ગણતરી છે, તેથી સંતાન-સૂચકમાં
  એક ઉમેરવું પડે. આખી
  $\fn{SubtreeSeq}(d)$નો સંકેત
  માર્ગના સંકેતોની મર્યાદા હોવો
  જરૂરી નથી; અહીં અગાઉ વ્યાખ્યાયિત
  $\fn{sequenceBound}$માંથી સ્પષ્ટ
  મર્યાદા લીધી છે. અને ખુલ્લી
  ધારણાનું ચિહ્નન~$0$ હોય ત્યારે
  પૂર્વજનાં શૂન્ય નિવૃત્તિ-ચિહ્નનો
  તેને નિવૃત્ત કરતો નથી; તે કિસ્સો
  અહીં અલગ રાખ્યો છે.}
\end{proof}

\begin{prop}
  \ollabel{prop:prf-prim-rec}
  $\Gamma$ એ !!{sentence}sનો આદિમ
  પુનરાવર્તી ગણ હોય એમ ધારો.
  ત્યારે $\Prf[\Gamma](x, y)$ આદિમ
  પુનરાવર્તી છે: તેનો અર્થ એ છે કે
  $x$ એ~$\Gamma$ની
  !!{undischarged} ધારણાઓમાંથી
  $!A$ની !!a{derivation}~$\delta$નો
  સંકેત છે અને $y$ એ~$!A$ની
  ગડેલ-સંખ્યા છે.
\end{prop}

\begin{proof}
  ``$y \in \Gamma$'' આદિમ પુનરાવર્તી
  વિધાન~$R_\Gamma(y)$ વડે આપવામાં
  આવ્યું હોય એમ ધારો. આપણે બતાવવું છે
  કે $\Prf[\Gamma](x, y)$ આદિમ
  પુનરાવર્તી છે. આ સંબંધ ત્યારે અને
  માત્ર ત્યારે સત્ય છે જ્યારે $y$ એ
  વાક્ય~$!A$ની ગડેલ-સંખ્યા હોય અને
  $x$ એ અંતિમ !!{formula}~$!A$
  તથા~$\Gamma$માંની બધી
  !!{undischarged} ધારણાઓ ધરાવતી
  સ્વાભાવિક નિગમનની
  !!{derivation}નો સંકેત હોય.

  \olref{prop:deriv} પ્રમાણે, $x$ એ
  સ્વાભાવિક નિગમનની યોગ્ય
  !!{derivation}~$\delta$ની ગડેલ-સંખ્યા
  હોય ત્યારે સત્ય થતો
  $\fn{Deriv}(x)$ ગુણધર્મ આદિમ
  પુનરાવર્તી છે. તેથી
  $\Prf[\Gamma](x, y)$ને આ પ્રમાણે
  વ્યાખ્યાયિત કરી શકીએ:
  \begin{align*}
    \Prf[\Gamma](x, y) \defiff {}
    & \fn{Deriv}(x) \land \fn{EndFmla}(x) = y \land {} \\
    & \bforall{z < x}{(\fn{OpenAssum}(z, x) \lif R_\Gamma(z))}.
  \end{align*}
\end{proof}

\end{document}
